\section{Introduction}

\emph{Chronic kidney disease} (\glossterm{ckd}{CKD}) is common and growing. In 2023 it affected around \textbf{14\% of adults} worldwide, about 788 million people, up from 378 million in 1990.\cite{gbd2023ckd} Its early stages usually have no symptoms, so many people do not know they have it until their kidneys are already damaged. This is one reason it was the \textbf{ninth leading cause of death} in 2023, with 1.5 million deaths.\cite{gbd2023ckd} \emph{Machine learning} can help here. Standard blood and urine tests, like \glossterm{creatinine}{creatinine}, \glossterm{haemoglobin}{haemoglobin} and \glossterm{albumin}{albumin}, are already part of normal check-ups, and a model trained on these values could flag likely CKD cases cheaply.

In this project, we use two separate datasets to investigate kidney health through different data modalities. The first is a tabular clinical dataset containing patient information and clinical measurements, which is used to predict CKD status. On this dataset, we compare four classical classifiers: \hyperref[sec:dt]{Decision Tree}, \hyperref[sec:knn]{kNN}, \hyperref[sec:nb]{Naive Bayes} and \hyperref[sec:rf]{Random Forest}. The second is the URI-CAD renal ultrasound dataset, containing 1,985 ultrasound images labelled as healthy (450 images) or pathological (1,535 images). We use these images to develop and evaluate \hyperref[sec:cnn]{Convolutional Neural Network (CNN)} models for binary image classification.\cite{uricad2024}

Although both datasets address the broader task of identifying kidney health conditions, they are independent datasets and are not directly linked or combined. The tabular experiment predicts CKD using clinical measurements, whereas the image experiment classifies ultrasound images according to their global diagnosis labels of healthy or pathological. Therefore, the two experiments provide separate approaches to investigating kidney health using clinical and imaging data.

The project investigates which classical model predicts CKD best; which features matter most and whether the engineered ones help; whether unsupervised \emph{clustering} groups patients in a way that matches the CKD diagnosis; and how effectively CNN models classify healthy and pathological kidney ultrasound images.

\vspace{2\baselineskip}